\documentclass[12pt,a4paper]{article}

% ==================================================
% 1. CONFIGURATION
% ==================================================

% Enable / disable title page
\newif\ifshowtitlepage
\showtitlepagetrue
% \showtitlepagefalse

% Enable / disable table of contents
\newif\ifshowtoc
\showtocfalse

% Document information
\newcommand{\worktitle}{Backlog Management Methodology}
\newcommand{\worksubtitle}{NAO Robot Exercise System}
\newcommand{\teamname}{NAOMI}
\newcommand{\subject}{NAO Robot Exercise System}

\newcommand{\university}{Slovak University of Technology in Bratislava}
\newcommand{\faculty}{Faculty of Informatics and Information Technologies}

\newcommand{\academicyear}{2026/2027}
\newcommand{\location}{Bratislava}
\newcommand{\documentyear}{2026/2027}

% Team leaders
\newcommand{\leaderone}{Ing. Jakub Perdek, PhD.}
\newcommand{\leadertwo}{doc. Ing. Fedor Lehocki, PhD.}

% Team members
\newcommand{\memberone}{Bc. Akhmadinurov Amal}
\newcommand{\membertwo}{Bc. Danilov Denis}
\newcommand{\memberthree}{Bc. Dmitriev Anton}
\newcommand{\memberfour}{Bc. Gainetdinov Ruslan}
\newcommand{\memberfive}{Bc. Kuzin Timofei}
\newcommand{\membersix}{Bc. Tikhoblazhenko Alena}
\newcommand{\memberseven}{Bc. Zimnitski Aliaksei}


% ==================================================
% 2. PACKAGES
% ==================================================

\usepackage[
    a4paper,
    top=3cm,
    bottom=2.7cm,
    left=2.8cm,
    right=2.8cm,
    headheight=45pt,
    headsep=20pt
]{geometry}

\usepackage[T1]{fontenc}
\usepackage[utf8]{inputenc}
\usepackage[english]{babel}
\usepackage{lmodern}

\usepackage{graphicx}
\graphicspath{{../../assets/}{assets/}}
\usepackage[table]{xcolor}
\usepackage{fancyhdr}
\usepackage{titlesec}
\usepackage{setspace}
\usepackage{lastpage}
\usepackage{array}
\usepackage{float}
\usepackage{caption}
\usepackage{booktabs}
\usepackage{tabularx}
\usepackage{lipsum}
\usepackage{csquotes}

% Hyperlinks
\usepackage[hidelinks]{hyperref}


% ==================================================
% 3. COLORS
% ==================================================

\definecolor{primary}{HTML}{000000}
\definecolor{secondary}{HTML}{000000}
\definecolor{muted}{HTML}{000000}
\definecolor{linecolor}{HTML}{000000}
\definecolor{naomiTableHeader}{HTML}{17445A}
\definecolor{naomiTableStripe}{HTML}{EDF7FA}
\definecolor{naomiTableRule}{HTML}{66BBD0}

% Ragged-right columns keep narrow table cells readable.
\newcolumntype{L}[1]{>{\raggedright\arraybackslash}p{#1}}
\newcolumntype{Y}{>{\raggedright\arraybackslash}X}


% ==================================================
% 4. TYPOGRAPHY
% ==================================================

\setstretch{1.2}

\setlength{\parindent}{0pt}
\setlength{\parskip}{7pt}

% Section
\titleformat{\section}
    {\Large\bfseries\color{primary}}
    {\thesection}
    {0.8em}
    {}

% Subsection
\titleformat{\subsection}
    {\large\bfseries\color{primary}}
    {\thesubsection}
    {0.8em}
    {}

% Subsubsection
\titleformat{\subsubsection}
    {\normalsize\bfseries\color{primary}}
    {\thesubsubsection}
    {0.8em}
    {}

\titlespacing*{\section}{0pt}{20pt}{10pt}
\titlespacing*{\subsection}{0pt}{14pt}{7pt}
\titlespacing*{\subsubsection}{0pt}{12pt}{5pt}

% Captions
\captionsetup{
    font=small,
    labelfont=bf,
    justification=centering
}


% ==================================================
% 5. HEADER AND FOOTER
% ==================================================

\pagestyle{fancy}
\fancyhf{}

% Header logo
\fancyhead[L]{
    \includegraphics[
        width=4cm,
        height=1.15cm,
        keepaspectratio
    ]{horizontal.png}
}

% Header information
\fancyhead[R]{
    \begin{tabular}{r}
        {\footnotesize\bfseries \worksubtitle}\\[3pt]
        {\scriptsize Academic Year \academicyear}
    \end{tabular}
}

% Header line
\renewcommand{\headrulewidth}{0.7pt}

% Footer left
\fancyfoot[L]{
    {\scriptsize
    \teamname\enspace|\enspace\subject}
}

% Footer right
\fancyfoot[R]{
    {\scriptsize
    \thepage\enspace/\enspace\pageref{LastPage}}
}

% Footer line
\renewcommand{\footrulewidth}{0.4pt}

% Plain pages
\fancypagestyle{plain}{
    \fancyhf{}

    \fancyhead[L]{
        \includegraphics[
            width=4cm,
            height=1.15cm,
            keepaspectratio
        ]{horizontal.png}
    }

    \fancyhead[R]{
        \begin{tabular}{r}
            {\footnotesize\bfseries \worksubtitle}\\[3pt]
            {\scriptsize Academic Year \academicyear}
        \end{tabular}
    }

    \fancyfoot[L]{
        {\scriptsize
        \teamname\enspace|\enspace\subject}
    }

    \fancyfoot[R]{
        {\scriptsize
        \thepage\enspace/\enspace\pageref{LastPage}}
    }

    \renewcommand{\headrulewidth}{0.7pt}
    \renewcommand{\footrulewidth}{0.4pt}
}




% ==================================================
% 6. TITLE PAGE
% ==================================================

\newcommand{\makemytitlepage}{
    \begin{titlepage}
        \thispagestyle{empty}
        \centering

        % ==========================================
        % UNIVERSITY LOGO
        % ==========================================

        \vspace*{0.3cm}

        \includegraphics[
            width=6cm,
            height=5cm,
            keepaspectratio
        ]{vertical.png}

        \vspace{0.5cm}

        {\large\bfseries
        \university\par}

        \vspace{0.2cm}

        {\normalsize
        \faculty\par}

        % ==========================================
        % PROJECT TITLE
        % ==========================================

        \vspace{1.7cm}

        {\small\MakeUppercase{\worksubtitle}\par}

        \vspace{0.5cm}

        \begin{minipage}{0.95\textwidth}
            \centering

            {\fontsize{25}{32}\selectfont\bfseries
            \worktitle\par}
        \end{minipage}

        % ==========================================
        % TEAM INFORMATION
        % ==========================================

        \vfill

        \begin{minipage}{0.9\textwidth}
            \centering
            \begin{minipage}[t]{0.46\linewidth}
                \centering
                {\normalsize\bfseries Team Leaders\par}
                \vspace{0.25cm}
                {\small
                \leaderone\par
                \leadertwo\par}
            \end{minipage}%
            \hfill
            \begin{minipage}[t]{0.48\linewidth}
                \centering
                {\normalsize\bfseries Team Members\par}
                \vspace{0.25cm}
                {\small
                \memberone\par
                \membertwo\par
                \memberthree\par
                \memberfour\par
                \memberfive\par
                \membersix\par
                \memberseven\par}
            \end{minipage}
        \end{minipage}

        % ==========================================
        % BOTTOM INFORMATION
        % ==========================================

        \vfill

        {\small
        \location\enspace\textbullet\enspace\academicyear\par}

    \end{titlepage}
}

% ==================================================
% 7. IMAGE PLACEHOLDER
% ==================================================

% Usage:
% \imageplaceholder{Text}{Caption}

\newcommand{\imageplaceholder}[2]{

    \begin{figure}[H]
        \centering

        \begingroup
        \setlength{\fboxsep}{0pt}

        \fbox{
            \begin{minipage}[c][5cm][c]{0.85\textwidth}

                \centering

                {\Large\bfseries IMAGE}

                \vspace{0.3cm}

                {\normalsize #1}

            \end{minipage}
        }

        \endgroup

        \caption{#2}

    \end{figure}
}


% ==================================================
% 8. REAL IMAGE COMMAND
% ==================================================

% Usage:
% \documentimage{path}{caption}{label}

\newcommand{\documentimage}[3]{

    \begin{figure}[htbp]
        \centering

        \includegraphics[
            width=0.85\textwidth,
            keepaspectratio
        ]{#1}

        \caption{#2}
        \label{#3}

    \end{figure}
}


% ==================================================
% 9. TABLES
% ==================================================

% Usage:
% \begin{naomitable}{Caption}{label}
%     \begin{tabularx}{\textwidth}{L{2.8cm} Y L{2.6cm}}
%         \rowcolor{naomiTableHeader}
%         \textcolor{white}{\bfseries Name} & ... \\
%         ...
%         \arrayrulecolor{naomiTableRule}\bottomrule[1pt]
%     \end{tabularx}
% \end{naomitable}

\newenvironment{naomitable}[2]{%
    \begin{table}[htbp]
        \centering
        \caption{#1}
        \label{#2}
        \begingroup
        \small
        \renewcommand{\arraystretch}{1.4}
        \setlength{\tabcolsep}{7pt}
        \rowcolors{2}{naomiTableStripe}{white}
}{%
        \endgroup
    \end{table}
}

% Usage:
% \tableplaceholder{Caption}{label}

\newcommand{\tableplaceholder}[2]{
    \begin{naomitable}{#1}{#2}
        \begin{tabularx}{\textwidth}{L{2.8cm} Y L{2.6cm}}
            \rowcolor{naomiTableHeader}
            \textcolor{white}{\bfseries Name} &
            \textcolor{white}{\bfseries Description} &
            \textcolor{white}{\bfseries Status} \\
            Item A & Sample description & Completed \\
            Item B & Sample description & In Progress \\
            Item C & Sample description & Planned \\
            \arrayrulecolor{naomiTableRule}\bottomrule[1pt]
        \end{tabularx}
    \end{naomitable}
}


% ==================================================
% 10. DOCUMENT
% ==================================================

\begin{document}

\ifshowtitlepage
    \makemytitlepage
\fi

\ifshowtoc
    \tableofcontents
    \clearpage
\fi

% ==================================================
% 1. PURPOSE
% ==================================================

\section{Purpose of the Document}

This document describes how the NAOMI team manages its project backlog
using Jira and a Scrumban approach. The aim is to make planned work
visible, keep responsibilities clear, and support regular delivery of
results throughout the academic year.

Scrumban combines sprint-based planning and review with the visual
workflow of a Kanban board. For NAOMI, this means agreeing on short-term
goals while allowing the team to monitor work in progress, identify
blockers, and adjust priorities when necessary.

% ==================================================
% 2. SCOPE
% ==================================================

\section{Scope and Responsibilities}

The methodology applies to all NAOMI team members working on analysis,
robot-control software, motion processing, testing, documentation, and
other project activities. Jira serves as the central place for recording
and following the status of these activities.

The Product Owner confirms product requirements and their priority
before related work enters a sprint.

The team lead coordinates sprint planning, monitors progress,
and helps resolve blockers. The deputy team lead supports these
activities and substitutes for the team lead when unavailable.

Team members estimate work, implement and review changes, validate
results, and keep their assigned Jira issues up to date.

The Product Owner is \textbf{Ing. Jakub Perdek, PhD.}
The team lead is Ruslan Gainetdinov, and the deputy team lead
is Alena Tikhoblazhenko.

% ==================================================
% 3. BACKLOG STRUCTURE
% ==================================================

\section{Backlog Organisation in Jira}

\subsection{Product Backlog}

The product backlog contains the features, improvements, research tasks,
experiments, and technical work that may be needed to reach the project's
goals. Items should have a concise title and a description of the expected
outcome. Larger activities may be represented as Epics and divided into
smaller Stories or Tasks, using the issue types available in the team's
Jira project.

The backlog is reviewed and reordered when requirements change or new
information becomes available. Higher-priority items should be
sufficiently clear before the team starts working on them.

\subsection{Sprint Backlog}

At the beginning of each sprint, the team selects a realistic set of
issues from the product backlog, taking into account available capacity,
priorities, and dependencies. The selected work is visible on the Jira
board. Unfinished issues are discussed during the sprint review or next
planning session and are not automatically treated as completed.

\subsection{Sprint Length and Meeting Rhythm}

Each sprint lasts \textbf{two weeks}.

Sprint planning takes place at the beginning of each sprint.
The team agrees on the sprint goal and selects work according
to priority, dependencies, and available capacity.

Regular team coordination meetings take place
\textbf{on Mondays or Tuesdays}. Members discuss progress,
blockers, and any necessary changes to the sprint backlog.

A sprint review and retrospective take place at the end of each
sprint. The team reviews completed results, discusses unfinished
work, and agrees on improvements for the next sprint.

Meetings with the supervisor take place
\textbf{on Fridays}.

\subsection{Requirements Before Sprint Admission}

New functionality and changes to product requirements are discussed
with the Product Owner before implementation tasks are created.

If a requirement first needs investigation, the team creates
an analysis or research task and presents the results to the
Product Owner. Implementation tasks are created after
the Product Owner approves the proposed requirements.

Technical and internal tasks, such as CI improvements,
methodology updates, bug fixes, and refactoring, do not require
Product Owner confirmation.

% ==================================================
% 4. SCRUMBAN WORKFLOW
% ==================================================

\section{Scrumban Workflow}

The NAOMI team uses a Jira board to manage tasks and
track their progress throughout the development process.
Our workflow consists of five stages, allowing team
members to monitor the current state of each issue.

\begin{naomitable}{Scrumban workflow states}{tab:workflow}
    \begin{tabularx}{\textwidth}{L{2.8cm} Y}
        \rowcolor{naomiTableHeader}
        \textcolor{white}{\bfseries Status} &
        \textcolor{white}{\bfseries Description} \\
        To Do & The task is planned but work has not started. \\
        In Progress & A team member is actively working on the task. \\
        In Review & The implementation is completed and waiting
        for review or feedback from other team members. \\
        Testing & A code task has passed review and is being tested before merge. \\
        Done & Applicable review and validation are complete; code changes have been merged. \\
        \arrayrulecolor{naomiTableRule}\bottomrule[1pt]
    \end{tabularx}
\end{naomitable}

Non-code tasks move directly from \textbf{In Review} to \textbf{Done}
after review and appropriate validation. Code tasks move from
\textbf{In Review} to \textbf{Testing}; testing is completed before the
pull request is merged. They move to \textbf{Done} after the merge.

Team members are responsible for keeping the status
of their assigned issues up to date in Jira.
Tasks should be moved between columns as soon as
their status changes, rather than only at the end
of a sprint.

If a task cannot be completed due to a blocker,
the responsible team member should document the
problem in Jira and inform the rest of the team.

Following the Scrumban approach, the team focuses
on completing tasks already in progress before
starting new ones. This helps avoid unnecessary
multitasking and improves the overall workflow.

\subsection{Work-in-Progress Limit}

The \textbf{In Progress} column has a team-wide WIP limit
of \textbf{10 issues}. This limit applies to the whole team,
not to each member individually.

When the limit is reached, members help complete or unblock
existing work before moving another issue to In Progress.
Blocked issues remain counted while they are in this column.


% ==================================================
% 5. REGULAR REVIEW
% ==================================================

\section{Reviewing Progress and Identifying Blockers}

The Jira board is checked during regular team meetings and whenever an
assignee makes meaningful progress. During a review, the team considers
whether the current sprint goal is still achievable and whether work is
moving through the board as expected.

\begin{itemize}
    \item \textbf{Priorities:} Are we working on the most important issues?
    Have new requirements changed the plan?
    \item \textbf{Progress:} Which issues have remained In Progress for too
    long? Is anything blocked or waiting for a decision?
    \item \textbf{Ownership:} Does each active issue have a clear assignee?
    Are issue descriptions and next steps understandable?
    \item \textbf{Delivery:} Are completed results tested, documented, and
    ready to demonstrate or integrate?
\end{itemize}

Issues that need support should be raised at the next team meeting or
communicated earlier through the team's agreed communication channel.
Critical blockers should not be left unresolved until the sprint ends.

Reviewers must review a pull request within \textbf{two working days}
after review is requested. The PR author requests review as soon
as the change is ready.

If the review is delayed, the author informs the team lead
so that the review can be reassigned or the blocker resolved.
A task awaiting required review is not considered Done
at the end of the sprint.

% ==================================================
% 6. ISSUE RULES
% ==================================================

\section{Rules for Managing Jira Issues}

\subsection{Creating and Updating Issues}

Each issue must have a concise title, a description of the
expected result, and relevant acceptance criteria.

Before a task enters a sprint, the following Jira fields
must be completed:

\begin{itemize}
    \item \textbf{Parent}: the parent issue in the backlog hierarchy.
    \item \textbf{Assignee}: the team member responsible for the task.
    \item \textbf{Sprint}: the sprint in which the task is planned.
    \item \textbf{Story points}: the agreed estimate of relative effort.
    \item \textbf{Priority}: the task's agreed priority level.
\end{itemize}

These requirements apply to sprint tasks.
Epics group related work and \textbf{are not added to sprints}.
Their child Stories or Tasks are selected for individual sprints.

Links to relevant code, designs, experiments, or documentation
are attached where useful.

Each member updates the status of their issues and records important
findings or blockers in Jira. If an issue is too broad or cannot be
completed within a reasonable period, it should be divided into smaller
pieces of work.

\subsection{Changes During a Sprint}

New tasks are first added to the product backlog. They may enter
the running sprint only if they block the sprint goal.
Otherwise, they remain in the backlog for a future sprint.

Any addition to the running sprint must be discussed with the team.
The reason for the addition and its impact on the planned work
must be recorded in the related Jira issue.

\subsection{Finishing Work}

An issue is moved to Done only when its agreed acceptance criteria are
satisfied. Depending on the type of work, this may involve a code review,
a successful test, an updated document, or a demonstration of the result.
The team uses its Definition of Done as the common reference
for deciding whether an item is finished.

% ==================================================
% 7. PLANNING AND ESTIMATION
% ==================================================

\section{Planning and Estimation}

The team has adopted the story point scale
\textbf{1, 2, 3, 5, 8}. Every task included in a sprint must
be estimated before it enters that sprint, including analysis,
design, testing, and documentation tasks.

Story points express relative effort, complexity, and uncertainty;
they do not represent a fixed number of working hours.
Issue types do not determine fixed estimates.

\begin{naomitable}{Adopted story point scale}{tab:estimates}
    \begin{tabularx}{\textwidth}{L{2.4cm} Y}
        \rowcolor{naomiTableHeader}
        \textcolor{white}{\bfseries Points} &
        \textcolor{white}{\bfseries Interpretation} \\
        1 & Very small, well-understood task \\
        2 & Small task with limited uncertainty \\
        3 & Moderate task requiring several steps \\
        5 & Larger task involving more effort or uncertainty \\
        8 & Complex activity that may need to be split \\
        \arrayrulecolor{naomiTableRule}\bottomrule[1pt]
    \end{tabularx}
\end{naomitable}

Tasks estimated at \textbf{8 points} should be considered
for splitting. Larger tasks must be split into smaller tasks
before sprint admission.
Estimates are agreed by the team during planning or backlog refinement.

Sprint capacity accounts for university deadlines and the
availability of individual members. Completed work, carried-over
tasks, and blockers are reviewed to improve subsequent planning.

\subsection{Priority Levels}

Priority determines the order in which work should be considered.
It does not replace the story point estimate.

\begin{naomitable}{Task priority levels}{tab:priorities}
    \begin{tabularx}{\textwidth}{L{2.4cm} Y}
        \rowcolor{naomiTableHeader}
        \textcolor{white}{\bfseries Priority} &
        \textcolor{white}{\bfseries Meaning} \\
        Highest & Critical work that blocks essential functionality
        or an imminent delivery and requires immediate attention. \\
        High & Important work needed for the sprint goal
        or to unblock other tasks. \\
        Medium & Normal planned work that contributes
        to project goals without immediate urgency. \\
        Low & Useful improvements that can wait
        while higher-priority work is completed. \\
        Lowest & Optional refinements with minimal
        impact on current project goals. \\
        \arrayrulecolor{naomiTableRule}\bottomrule[1pt]
    \end{tabularx}
\end{naomitable}

% ==================================================
% 8. CONTINUOUS IMPROVEMENT
% ==================================================

\section{Continuous Improvement}

At the end of each sprint, the team briefly reviews what worked well,
what slowed progress, and which changes would help in the next sprint.
Possible improvements include clearer Jira issue descriptions, more
consistent board updates, smaller tasks, and adjustments to meeting or
review practices.

The purpose of this methodology is not to introduce unnecessary
administration. Jira should give the team a reliable overview of its
work while leaving enough time for actual development and experimentation.

\end{document}
